\subsection{积分次序的交换}

设 \(I = [a, b]\) 与 \(A = [c, d]\) 是 \(R\) 中的两个紧区间；设 \(f\) 是 \(I \times A\) 上的连续函数，取值于完备赋范空间 \(E\)（定义在 \(R\) 上）中；由 II, p.~69 的命题 2，\(\int_{a}^{b} f(x, \alpha) dx\) 是 \(\alpha\) 在 \(A\) 上的连续函数；为简单起见，它的积分 \(\int_{c}^{d} \left( \int_{a}^{b} f(x, \alpha) dx \right) d\alpha\) 也记作 \(\int_{c}^{d} d\alpha \int_{a}^{b} f(x, \alpha) dx\)。

命题 8. 若 \(\mathbf{f}\) 在 \(\mathrm{I} \times \mathrm{A}\) 上连续，则有

\[
\int_ {c} ^ {d} d \alpha \int_ {a} ^ {b} \mathbf {f} (x, \alpha) d x = \int_ {a} ^ {b} d x \int_ {c} ^ {d} \mathbf {f} (x, \alpha) d \alpha\tag{15}
\]

（“交换积分次序的公式”）。

我们将证明，对所有 \(y \in A\)，有

\[
\int_ {c} ^ {y} d \alpha \int_ {a} ^ {b} \mathbf {f} (x, \alpha) d x = \int_ {a} ^ {b} d x \int_ {c} ^ {y} \mathbf {f} (x, \alpha) d \alpha .\tag{16}
\]

由于 (16) 的两端都是 \(y\) 的函数，且在 \(y = c\) 时相等，故只需证明它们在 \(]c\) , \(d[\) 上可微且其导数在这个区间的每一点处相等。若令 \(\mathbf{g}(\alpha) = \int_{a}^{b} \mathbf{f}(x, \alpha) dx\) 与 \(\mathbf{h}(x, y) = \int_{c}^{y} \mathbf{f}(x, \alpha) dx\)，则关系式 (16) 可以写成

\[
\int_ {c} ^ {y} \mathbf {g} (\alpha) d \alpha = \int_ {a} ^ {b} \mathbf {h} (x, y) d x.
\]

而第一项对 \(y\) 的导数是 \(\mathbf{g}(y)\)，第二项的导数是 \(\int_{a}^{b} \mathbf{h}_{y}'(x, y) dx\)（由 II, p.74, 推论），因为 \(\mathbf{h}_{y}'(x, y) = \mathbf{f}(x, y)\) 在 \(I \times A\) 上连续；这样得到的两个表达式是相同的。

现在设 A = {[}c, d{]} 是 R 中的紧区间，I 是 R 中的任意区间；设 f 是 I × A 上的连续函数，取值于 E 中，使得积分 \(\mathbf{g}(\alpha) = \int_{\mathrm{I}} \mathbf{f}(t, \alpha) dt\) 对所有 \(\alpha \in A\) 收敛；即使 \(\mathbf{g}(\alpha)\) 在 A 上连续，也未必总能在积分 \(\int_{c}^{d} d\alpha \int_{I} \mathbf{f}(t, \alpha) dt\) 中交换积分次序，因为积分 \(\int_{I} dt \int_{c}^{d} \mathbf{f}(t, \alpha) d\alpha\) 可能不存在，或者可能与积分 \(\int_{c}^{d} d\alpha \int_{I} \mathbf{f}(t, \alpha) dt\) 不同（参见 II, p.~87, 习题 7）。然而有下列结果：

命题 9. 若函数 \(\mathbf{f}\) 在 \(\mathrm{I} \times \mathrm{A}\) 上连续，且积分 \(\int_{\mathrm{I}} \mathbf{f}(t, \alpha) dt\) 在 \(\mathrm{A}\) 上一致收敛，则积分 \(\int_{\mathrm{I}} dt \int_{c}^{d} \mathbf{f}(t, \alpha) d\alpha\) 收敛，且有

\[
\int_ {c} ^ {d} d \alpha \int_ {\mathrm{I}} \mathbf {f} (t, \alpha) d t = \int_ {\mathrm{I}} d t \int_ {c} ^ {d} \mathbf {f} (t, \alpha) d \alpha .\tag{17}
\]

对含于 I 的每个紧区间 J，令 \(\mathbf{u}_{\mathrm{J}}(\alpha) = \int_{\mathrm{J}} \mathbf{f}(t, \alpha) dt\)。这个假设蕴含：关于滤子 \(\Phi\)（有向集 \(\Re(\mathrm{I})\) 的截口滤子），连续函数 \(u_{J}\) 在 A 上一致收敛于 \(\int_{I} f(t, \alpha) dt\)；于是（II, p.~68, 命题 1），\(\int_{c}^{d} d\alpha \int_{J} f(t, \alpha) dt\) 有极限 \(\int_{c}^{d} d\alpha \int_{I} f(t, \alpha) dt\)（关于 \(\Phi\)）；但由命题 8（II, p.~77），有

\[
\int_ {c} ^ {d} d \alpha \int_ {\mathrm{J}} \mathbf {f} (t, \alpha) d t = \int_ {\mathrm{J}} d t \int_ {c} ^ {d} \mathbf {f} (t, \alpha) d \alpha .\tag{18}
\]

前面的结果于是表明积分 \(\int_{I} dt \int_{c}^{d} \mathbf{f}(t, \alpha) d\alpha\) 收敛，而且在关系式 (18) 中关于 \(\Phi\) 取极限，便得到 (17)。

\setcounter{exercisegroup}{0}
\refstepcounter{exercisegroup}
